\chapter[Band 42: Endliche Gruppen]{Band 42 auf einen Blick}
\noindent\textbf{Endliche Gruppen}\par\medskip
Der Band besitzt einen eigenen Begriff der endlichen Gruppe.
Sein ausgearbeiteter Schwerpunkt ist die Rekonstruktion endlicher
Halbgruppen mit einem \emph{Gruppenquadrat}. Benötigt werden insbesondere
Potenzhalbgruppen, Inflationen, Faserbijektionen und endliche
Potenzmengenkardinalitäten.

\begin{BandUebersicht}
\textbf{Eigene Struktur: endliche Gruppe}
\UebFormel{\begin{aligned}
&\mathsf{FinGrp}(G,\star,e)\ \leftrightarrow\\
&\GroupAlg G\star e\land\Finite G.
\end{aligned}}Endlichkeit ist hier ausdrücklich Bestandteil des Strukturbegriffs.
\UebQuellen{\FormulaRefAuto{FiniteGroupDef}[def]}
&
\textbf{Einordnung}
\UebFormel{\begin{aligned}
&\mathsf{FinGrp}(G,\star,e)\ \vdash\\
&\qquad\FiniteSemigroupAlg G\star.
\end{aligned}}Insbesondere sind die Sätze über endliche Halbgruppen auf den
zugrunde liegenden Halbgruppenträger anwendbar.
\UebQuellen{\FormulaRefAuto{FiniteGroupIsFiniteSemigroup}[thm]} \\
\addlinespace[.8ex]

\textbf{Gruppenquadrat und Potenzisomorphismus}
Für nichtleere Halbgruppen \((S,\star)\), \((T,\diamond)\) setze\UebFormel{\begin{aligned}
&G=S^2,\quad H=T^2,\\
&Q_S=\PowNE S,\\
&Q_T=\PowNE T.
\end{aligned}}Es sei \(G\) eine Gruppe mit Eins \(e\) und \(F\) ein
Isomorphismus von \((Q_S,\star_{\mathcal P})\) nach
\((Q_T,\diamond_{\mathcal P})\).
\UebQuellen{\FormulaRefAuto{SemigroupSquareDef}[def];
\FormulaRefAuto{SemigroupIsoDef}[def];
\FormulaRefAuto{GroupDef}[def]}
&
\textbf{Der Gruppenkern wird rekonstruiert}
Dann existieren \(u\in H\) und ein Gruppenisomorphismus
\(\varphi:(G,\star,e)\cong(H,\diamond,u)\) mit\UebFormel{\begin{aligned}
&g\in G\ \vdash\
 F(\{g\})=\{\varphi(g)\}.
\end{aligned}}Für diesen Kernschritt ist noch keine Endlichkeit nötig.
\UebQuellen{\FormulaRefAuto{PowerGroupSquareCoreReconstruction}[thm]} \\
\addlinespace[.8ex]

\textbf{Grundfasern der Inflation}
Sind nun beide Halbgruppen endlich und sind \(u,\varphi\) wie
zuvor gewählt, so seien\UebFormel{\begin{aligned}
&r(x)=e\star x,\qquad s(y)=u\diamond y,\\
&C_g=\{x\in S\mid r(x)=g\},\\
&D_h=\{y\in T\mid s(y)=h\}.
\end{aligned}}
\UebQuellen{\FormulaRefAuto{MonoidSquareInflationDef}[def]}
&
\textbf{Endliche Fasern sind gleichmächtig}
Unter allen Voraussetzungen des vorigen Schritts und
\(\Finite S,\Finite T\) gilt:\UebFormel{\begin{aligned}
&\forall g\in G\ \exists f_g:
 C_g\bij D_{\varphi(g)}.
\end{aligned}}Der Potenzisomorphismus liefert zunächst Potenzfaserbijektionen;
die Endlichkeit erlaubt den Rückschluss auf die Grundfasern.
\UebQuellen{\FormulaRefAuto{FiniteGroupSquareFiberBijections}[thm]} \\
\addlinespace[.8ex]

\textbf{Zusammenfügen über dem Gruppenquadrat}
Das Quadrat ist eine Unterhalbgruppe; die Rückführungen
\(r:S\to G\), \(s:T\to H\) machen \(S,T\) zu Inflationen.
Die vorige Faserinformation wird mit dem Gruppenisomorphismus verklebt.
\UebQuellen{\FormulaRefAuto{MonoidSquareInflationProperties}[thm];
\FormulaRefAuto{FiniteGroupSquareFiberBijections}[thm]}
&
\textbf{Rekonstruktionssatz}
Für endliche Halbgruppen \(S,T\) gilt:\UebFormel{\begin{aligned}
&S\ne\varnothing,\quad
 \GroupAlg{S^2}\star e,\\
&(\PowNE S,\star_{\mathcal P})
 \cong(\PowNE T,\diamond_{\mathcal P})\\
&\qquad\vdash (S,\star)\cong(T,\diamond).
\end{aligned}}Die Nichtleerheit von \(T\) folgt bereits aus diesen Voraussetzungen.
\UebQuellen{\FormulaRefAuto{FiniteGroupSquarePowerSemigroupReconstructionOneSidedNonempty}[thm]} \\
\addlinespace[.8ex]
\end{BandUebersicht}

